\section{Коррекция ошибок --- это топология}
\label{sec:ecc}

\subsection{Проводка --- это код}

В обычной машине отказоустойчивость --- отдельный слой: ECC-DRAM выделяет
проверочные биты (накладные $+12{,}5\%$ хранения для распространённого
SEC-DED-кода $[72,64]$) и декодер синдрома; квантовая машина тратит порядка
$10^3$ физических кубитов на логический на surface-коде. Накладные существуют,
потому что код --- структура, \emph{добавленная к} вычислению. \TALOS{} не
добавляет ничего, потому что интерконнект уже \emph{есть} код дистанции 3.

\begin{claimbox}[Локализация отказа с нулевыми накладными \status{Т}]
Инцидентность Фано --- проверочная матрица кода Хэмминга $[n,k,d]=[7,4,3]$:
одиночные ошибки исправляются, двойные детектируются, а один деградировавший
сектор \emph{локализуется} тремя синдромными проверками --- пирамида
диагностики $21\to 7\to 3\to 1$ (двадцать одно пар-отношение коллапсирует к семи
осям, считываемым тремя синдромными битами, дающими один вердикт отказа). Три
проверки чётности --- три прямые через отказавшую точку. Нет хранения, нет
массива декодера и нет рассогласования код-против-вычисления: фабрика
\S\ref{sec:fabric} \emph{есть} проверочная матрица.
\end{claimbox}

Синдромное отображение точно и было машинно-проверено: для каждого из семи
одно-секторных паттернов отказа трёхбитный синдром, прочитанный как двоичный
индекс, возвращает позицию отказа (Приложение~\ref{app:fano}). Поскольку код ---
код Хэмминга, его гарантии учебниковые ($d=3$: исправить 1, детектировать 2),
унаследованные при нулевой предельной цене.

\paragraph{Инвариант прямоугольника: второй, более глубокий слой чётности (T-281/T-282).}
Слой Хэмминга выше сторожит \emph{оси} (какой сектор сбоит). Генезисные
результаты корпуса добавляют стража уровнем ниже: \emph{знаковую прошивку}
самой ткани. Жизнеспособность знаковой таблицы Фано --- «никакие два состояния
не могут аннигилировать» --- доказуемо \emph{линейное} условие: 84 правила
прямоугольника $F(p,r)F(q,s)F(p,s)F(q,r) = -1$ по XOR-классам, чьё
исчерпывающее множество решений --- ровно одна калибровочная орбита октонионов
(корпусная T-281). Бит-флип в хранимой знаковой таблице потому либо попадает
внутрь 16-элементной калибровочной орбиты (безвредная переметка), либо
нарушает хотя бы один прямоугольник --- \emph{84-битный XOR-самотест
прошивки}, дешёвый как проход чётности, сертифицирующий, что закон умножения
машины всё ещё жизнеспособный. И та же алгебра объясняет, почему никакой патч
прошивки не «расширит» ткань: на $\Omega^4$ соответствующая система
\emph{несовместна} (960 ограничений, ранг 214 --- корпусная T-282), так что
четвёртое зеркало --- не инженерный вызов, а противоречие. Хэмминг проверяет
состояние; прямоугольный слой проверяет \emph{закон}.

\subsection{Новый примитив: вычисление есть самодиагностика}

Есть следствие, не имеющее аналога в слоистой машине, и это один из подлинно
новых архитектурных фактов, которые обнажает \TALOS{}.

\begin{principle}[Тождество вычисление-диагностика]\label{prin:selfdiag}
Всякий тик и вычисляет, и эмитит собственный синдром. Кубическая часть
$\LindOmega$ порождается произведениями по прямым (закон замыкания): здоровая
тройка когерентностей \emph{замыкается} на своей прямой, а когерентность,
которая замкнуться не может, есть, по той же арифметике, ненулевой синдром.
Детекция ошибок поэтому не отдельный проход по состоянию --- это побочный
продукт его продвижения. Машина не может считать, не проверяя себя
одновременно.
\end{principle}

\noindent\emph{Почему это ново.} В каждой прежней архитектуре проверка --- лишняя
работа: скраббер метёт DRAM, декодер бежит после чтения, синдром извлекается
выделенными тактами. Слоистый проект \emph{вынужден} отделять счёт от проверки,
потому что код и вычисление --- разные структуры. Под коллапсом они одна
структура, поэтому проверка бесплатна в строгом смысле, что её удаление
потребовало бы удаления вычисления. \TALOS{} непрерывно самодиагностируется без
добавленной цены --- петля самодиагностики DIAKRISIS программного слоя \SYNARC{}
на аппаратном уровне есть просто чтение синдрома, который тик уже произвёл.

\subsection{Отказоустойчивая композиция}

Отказоустойчивость переживает композицию. Поскольку любые две прямые
проективной плоскости встречаются, потеря любого одного сектора оставляет все
$\binom{6}{2}=15$ выживших пар по-прежнему делящими шину (проверено для каждого
удалённого сектора), поэтому холон деградирует плавно, а не разбивается. На
уровне федерации грамматика Голея трёх организмов $[23,12,7]$ расширяет тот же
принцип на многоотказовую устойчивость по составленным холонам (приложение
$\Omega$-организма \SYNARC{}); \TALOS{} наследует защиту дистанции 3 на уровне
одного холона, а федерация наследует защиту дистанции 7 над ним, оба из того же
проективного семейства.

\begin{table}[H]
\centering\small
\caption{Отказоустойчивость: привинченные коды против кода-как-топологии.}
\label{tab:ecc}
\begin{tabular}{@{}p{3.6cm}p{3.4cm}p{4.7cm}@{}}
\toprule
\textbf{Схема} & \textbf{Накладные} & \textbf{Связь с вычислением} \\
\midrule
ECC-DRAM (SEC-DED) & $+12{,}5\%$ бит + декодер & отдельно; скраббинг --- лишние такты \\
Surface-код (кванты) & $\sim\!10^3$ физ./логич. & отдельно; декодирует против декогеренции \\
\TALOS{} (Фано $=$ Хэмминг) & \textbf{ноль} лишнего железа & \textbf{тождественно}: проводка есть проверка; всякий тик самодиагностируется (Принцип~\ref{prin:selfdiag}) \\
\bottomrule
\end{tabular}
\end{table}
